\section{Odbočka k~pravděpodobnosti}\label{sec:krypto-pravdepodobnost}

V~historických šifrách se často ptáme, kolik klíčů musí záškodník vyzkoušet. Moderní kryptografie však používá také náhodně generované klíče a~algoritmy, které při běhu provádějí náhodné volby. Bezpečnost potom nepopisujeme pouze větou „útok je těžký“, ale pravděpodobností, s~jakou může uspět. Proto si nejprve vyložíme základní pojmy potřebné v~dalších sekcích.

\subsection{Náhodný pokus, výsledky a~jevy}

Náhodný pokus je postup, jehož konkrétní výsledek předem neznáme, přestože známe množinu všech možných výsledků. Příkladem je hod hrací kostkou, hod mincí nebo rovnoměrný výběr bitového řetězce.

\begin{definition}
    \emph{Výběrový prostor} $\Omega$ je množina všech možných výsledků náhodného pokusu. \emph{Náhodný jev} je libovolná podmnožina $A\subseteq\Omega$. Jev nastane právě tehdy, když výsledek pokusu patří do množiny $A$.
\end{definition}

Při hodu šestistěnnou kostkou máme
\[
    \Omega=\set{1,2,3,4,5,6}.
\]
Jev „padne sudé číslo“ odpovídá množině $A=\set{2,4,6}$ a~jev „padne číslo větší než čtyři“ množině $B=\set{5,6}$. Průnik $A\cap B=\set{6}$ nastane, když nastanou oba jevy současně. Sjednocení $A\cup B=\set{2,4,5,6}$ nastane, když nastane alespoň jeden z~nich.

\begin{figure}[ht]
    \centering
    \input{06-kryptografie/images/ch06-pravdepodobnost-prostor.tex}
    \caption{Výběrový prostor hodu kostkou a~dva náhodné jevy.}
    \label{fig:krypto-pravdepodobnost-prostor}
\end{figure}

\begin{definition}
    Jsou-li všechny výsledky konečného výběrového prostoru $\Omega$ stejně pravděpodobné, definujeme pravděpodobnost jevu $A\subseteq\Omega$ vztahem
    \[
        \Pr(A)=\frac{\sizeof{A}}{\sizeof{\Omega}}.
    \]
\end{definition}

Pro sudé číslo na kostce tedy dostáváme
\[
    \Pr(A)=\frac{3}{6}=\frac{1}{2}.
\]
Pravděpodobnost nemožného jevu je $\Pr(\varnothing)=0$ a~pravděpodobnost jistého jevu je $\Pr(\Omega)=1$. Často ji zapisujeme také v~procentech, například $\frac{1}{2}=50\,\%$.

Jev, že $A$ nenastane, značíme $A^\complement=\Omega\setminus A$. Platí
\[
    \Pr(A^\complement)=1-\Pr(A).
\]
Tento jednoduchý vztah bude později velmi užitečný u~narozeninového paradoxu: místo přímého počítání pravděpodobnosti, že vznikne alespoň jedna shoda, nejprve spočítáme pravděpodobnost, že nevznikne žádná.

\paragraph{Více kroků náhodného pokusu.}

Při dvou hodech mincí je potřeba rozlišovat pořadí výsledků:
\[
    \Omega=\set{00,01,10,11},
\]
kde například $01$ znamená, že v~prvním hodu padla nula a~ve druhém jednička. Jsou-li oba hody nezávislé a~obě hodnoty stejně pravděpodobné, má každý výsledek pravděpodobnost $\frac{1}{4}$.

Obecně existuje právě $2^n$ bitových řetězců délky $n$. Vybereme-li řetězec $x$ rovnoměrně náhodně z~$\set{0,1}^{n}$, potom pro každý předem zvolený řetězec $a\in\set{0,1}^{n}$ platí
\[
    \Pr_{x\in\set{0,1}^{n}}[x=a]=\frac{1}{2^n}.
\]
Dolní index zde říká, odkud a~jak náhodnou hodnotu vybíráme; výraz v~hranatých závorkách popisuje jev, jehož pravděpodobnost měříme.

\begin{example}
    Vyberme rovnoměrně náhodně $x\in\set{0,1}^{4}$. Celkem existuje $2^4=16$ možností. Jev, že $x$ začíná jedničkou, obsahuje osm řetězců, a~proto má pravděpodobnost $\frac{8}{16}=\frac{1}{2}$. Jev $x=1011$ obsahuje jediný výsledek a~má pravděpodobnost $\frac{1}{16}$.
\end{example}

\subsection{Podmíněná pravděpodobnost a~nezávislost}

Někdy už o~výsledku něco víme a~ptáme se, jak tato informace změní pravděpodobnost jiného jevu. Jestliže například víme, že na kostce padlo sudé číslo, zůstávají možné jen výsledky $\set{2,4,6}$.

\begin{definition}
    Nechť $A,B\subseteq\Omega$ a~$\Pr(B)>0$. \emph{Podmíněnou pravděpodobnost} jevu $A$ za podmínky $B$ definujeme jako
    \[
        \Pr(A\mid B)=\frac{\Pr(A\cap B)}{\Pr(B)}.
    \]
\end{definition}

Pro jevy $A=\set{2,4,6}$ a~$B=\set{4,5,6}$ při hodu kostkou je $A\cap B=\set{4,6}$. Proto
\[
    \Pr(A\mid B)
    =
    \frac{\Pr(\set{4,6})}{\Pr(\set{4,5,6})}
    =
    \frac{2/6}{3/6}
    =
    \frac{2}{3}.
\]

\begin{definition}
    Jevy $A$ a~$B$ jsou \emph{nezávislé}, jestliže informace o~nastání jednoho nemění pravděpodobnost druhého, tedy
    \[
        \Pr(A\mid B)=\Pr(A).
    \]
    Je-li $\Pr(B)>0$, je tato podmínka ekvivalentní rovnosti
    \[
        \Pr(A\cap B)=\Pr(A)\Pr(B).
    \]
\end{definition}

\begin{figure}[ht]
    \centering
    \begin{tikzpicture}[
    font=\small,
    level 1/.style={sibling distance=6cm, level distance=1.4cm},
    level 2/.style={sibling distance=2.7cm, level distance=1.4cm},
    every node/.style={align=center},
    result/.style={rounded corners=1mm, draw=blue!60!black, fill=blue!10,
                   minimum width=1.2cm, minimum height=6mm},
    edge from parent/.style={draw=black!70, -{Latex}}
]
    \node {\(\Omega\)}
        child {node {\(0\)}
            child {node[result] {\(00\)\\\(\frac14\)}
                edge from parent node[left] {\(\frac12\)}}
            child {node[result] {\(01\)\\\(\frac14\)}
                edge from parent node[right] {\(\frac12\)}}
            edge from parent node[left] {\(\frac12\)}
        }
        child {node {\(1\)}
            child {node[result] {\(10\)\\\(\frac14\)}
                edge from parent node[left] {\(\frac12\)}}
            child {node[result] {\(11\)\\\(\frac14\)}
                edge from parent node[right] {\(\frac12\)}}
            edge from parent node[right] {\(\frac12\)}
        };
\end{tikzpicture}

    \caption{Dva nezávislé hody mincí a~pravděpodobnosti výsledných řetězců.}
    \label{fig:krypto-pravdepodobnost-strom}
\end{figure}

V~kryptografii je nezávislost klíče na zprávě zásadní. Pokud by volba klíče závisela na obsahu zprávy, mohl by způsob jeho generování sám prozrazovat informace. U~jednorázové šifry budeme dokonce požadovat, aby znalost šifrového textu nezměnila pravděpodobnost žádné původní zprávy:
\[
    \Pr(M=m\mid C=c)=\Pr(M=m).
\]

\paragraph{Náhodné veličiny.}

\begin{definition}
    \emph{Náhodná veličina} přiřazuje každému výsledku náhodného pokusu nějakou hodnotu. Formálně jde o~zobrazení $\mapping{X}{\Omega}{S}$, kde $S$ je množina možných hodnot veličiny.
\end{definition}

Písmena $M$, $K$ a~$C$ budeme používat pro náhodné veličiny představující postupně zprávu, klíč a~šifrový text. Zápis $M=m$ označuje jev, že náhodná veličina $M$ nabyla konkrétní hodnoty $m$. Rozdíl mezi $M$ a~$m$ je tedy podobný rozdílu mezi proměnnou a~její právě zvolenou hodnotou.

\notebox{Pravděpodobnostní algoritmus může při běhu používat vlastní náhodné bity. Pokud píšeme pravděpodobnost jeho úspěchu, počítáme nejen s~náhodným vstupem, ale také s~těmito vnitřními volbami. Stejný algoritmus proto může na stejném vstupu v~různých bězích postupovat odlišně.}
